\documentclass{situationist-preprint}

\title{Clifford, James and the Ethics of Belief}
\msid{EPI-23-a}
\msfamily{EPI-23}
\msdiscipline{Epistemology}
\mstype{Principal manuscript}
\mskeywords{ethics of belief; Clifford; James; will to believe; evidentialism; uniqueness; acceptance; stakes}

\begin{document}
\maketitle

\begin{abstract}
Clifford held that it is wrong always, everywhere and for anyone to believe on insufficient evidence; James replied that where a genuine option cannot be decided on intellectual grounds, we have the right to believe. We argue that each principle holds in its own domain. Clifford's principle holds where the stakes are high and evidence is available; James's right to believe applies where neither condition holds. The key is to identify the stakes that matter to Clifford's argument as the avoidable stakes of error: the harm that acting on a false belief risks and that withholding judgment and inquiring further would avoid. Where these are high and further evidence can be had, belief on insufficient evidence results from a culpable failure of inquiry, which is the wrong Clifford described. Where the option is forced, withholding avoids no risk, and where evidence cannot decide, there is no inquiry to fail at; Clifford's grounds then lapse. We reconcile James's right with the uniqueness of rational credence: where the evidence cannot decide, the unique rational credence is intermediate, and what James licenses is not a different credence but a commitment to act on the hypothesis.
\end{abstract}

\section{The problem}

\citet{Clifford1877} opened his essay with a shipowner who sends an old and often-repaired emigrant ship to sea. He has doubts about her seaworthiness, but he puts them aside, reflecting on her past voyages and on Providence, and comes to believe sincerely that she is safe. The ship goes down with all on board. Clifford's verdict was that the shipowner was guilty of the deaths, since he had no right to believe on such evidence, and that his guilt would have been no less had the ship arrived safely. He drew the general conclusion: \enquote{it is wrong always, everywhere, and for anyone, to believe anything upon insufficient evidence}.

\citet{James1896} replied that Clifford's rule, applied universally, would itself be irrational. An option between hypotheses is \emph{genuine} if it is live (each hypothesis is a real possibility for the agent), forced (there is no third course, since declining to choose has the effect of choosing one side) and momentous (the opportunity is unique, much depends on it, and the decision is irreversible). Where a genuine option cannot by its nature be decided on intellectual grounds, our passional nature may lawfully decide it, since to refuse to decide is itself a passional decision, which prefers the risk of losing the truth to the risk of error. James held that the religious hypothesis presents such an option, and that some truths, such as whether another person likes me, may depend on my prior faith in them.

The debate continues in contemporary form. \citet{vanInwagen1996} defended Clifford's principle and found it hard to apply consistently to philosophical and political beliefs; evidentialists separate the epistemic question, which attitude fits the evidence, from the moral question, what one ought to do \citep{ConeeFeldman1985}; and \citet{Pace2011} has argued that James's argument is best read as a case for the moral encroachment of practical considerations on justification. We argue that each of Clifford's and James's principles is correct within a domain fixed by two conditions, and we reconcile James's right to believe with the uniqueness of rational credence defended in EPI-17.

\section{Framework}

\begin{definition}[Sufficient evidence]\label{d:sufficient}
An agent's evidence is \emph{sufficient} for $p$ if the rational credence in $p$ on that evidence meets the threshold for belief.
\end{definition}

By EPI-17, for each body of evidence there is a unique rational credence in each proposition, so Definition~\ref{d:sufficient} is determinate once the threshold is fixed.

\begin{definition}[Stakes of error]\label{d:stakes}
The \emph{stakes of error} on $p$ are \emph{high} if acting on $p$, were $p$ false, would risk serious harm, to the agent or to others, that the agent could avoid by withholding judgment and not acting on $p$ until the question is better settled.
\end{definition}

\begin{definition}[Availability of evidence]\label{d:available}
Evidence on $p$ is \emph{available} if the agent can obtain further evidence bearing on $p$, at a cost proportionate to the stakes, before she must act.
\end{definition}

\begin{definition}[Commitment]\label{d:commit}
A \emph{commitment} to $p$ is a policy of acting on $p$, taking $p$ as a premise in practical deliberation and directing attention toward it, without disbelieving $p$.
\end{definition}

Definition~\ref{d:commit} is the notion of practical commitment used in REL-06-a and REL-17-a, and it is close to acceptance in the sense of \citet{Cohen1992}. Unlike belief, commitment is voluntary, since by EPI-11 belief is not under direct voluntary control while action and attention are.

\begin{principle}[Fit]\label{p:fit}
The epistemically permitted credence in $p$ is the unique credence that is rational on the agent's evidence; full belief is epistemically permitted only if the evidence is sufficient.
\end{principle}

\begin{principle}[Practical permission]\label{p:practical}
A commitment to $p$ is practically permitted if its expected value, by the rational credence, is at least that of the alternatives.
\end{principle}

\begin{definition}[Clifford's and James's principles]\label{d:principles}
\emph{Clifford's principle} is the claim that believing $p$ on insufficient evidence, and acting on that belief, is a wrong for which the agent is answerable. \emph{James's right to believe} is the claim that an agent may commit to $p$ on insufficient evidence without doing wrong.
\end{definition}

Definition~\ref{d:principles} states James's right as a right to commit. We argue in Section~\ref{s:unique} that this is the only form of the right compatible with uniqueness and with involuntarism, and that it is the form James's argument supports.

The distinction between the stakes of error and the importance of an option is central to what follows. James's momentous options are options on which much depends. But if an option is also forced, withholding judgment does not avoid the risk, since withholding has the same practical effect as deciding against the hypothesis: the agent who suspends judgment on the religious hypothesis forgoes the goods of religious life exactly as the unbeliever does. In a forced option, then, the stakes of error, in the sense of Definition~\ref{d:stakes}, are not high, however momentous the option. Clifford's shipowner faced no forced option: he could have delayed the voyage, had the ship overhauled, or sent her out with fewer passengers, and each of these would have avoided the harm his error risked.

\section{The argument}

\begin{theorem}[Clifford's domain]\label{t:clifford}
Where the stakes of error on $p$ are high and evidence on $p$ is available, Clifford's principle holds: an agent who believes $p$ on insufficient evidence and acts on it does wrong.
\end{theorem}

\begin{proof}
Since evidence is available, the agent could, by inquiry, have obtained evidence that either was sufficient for $p$ or showed that $p$ was not to be relied on. Since the stakes are high, withholding and inquiring would have avoided a risk of serious harm, at a cost proportionate to it. By EPI-11, belief follows attention, so a belief on insufficient evidence formed in these circumstances results from what the agent did or failed to do with her attention: she did not inquire, or she attended selectively to the evidence that favoured $p$. Either is an act or omission under her voluntary control, which exposed others or herself to an avoidable risk of serious harm for no proportionate benefit. She is answerable for it, and for the belief and the action that resulted.
\end{proof}

The proof locates Clifford's wrong where his own description places it. The shipowner did not decide to believe the ship was safe; he \enquote{put his doubts aside}, that is, he withdrew his attention from the evidence that bore against her soundness and dwelt on the considerations that comforted him. His evidence did not support the soundness of the ship, and on the reading of REL-17-a it opposed it. The proof also explains Clifford's claim that the shipowner would have been guilty even had the ship arrived: the wrong lies in exposing others to an avoidable risk, and it is complete before the outcome is known.

\begin{theorem}[James's domain]\label{t:james}
Where the stakes of error on $p$ are not high and evidence on $p$ is not available, James's right to believe applies: an agent who commits to $p$ on insufficient evidence does no wrong, provided the commitment is practically permitted.
\end{theorem}

\begin{proof}
Clifford's grounds, as the proof of Theorem~\ref{t:clifford} shows, are two: that a better-supported verdict could have been reached by inquiry, and that withholding would have avoided a serious risk. Where evidence is not available, there is no inquiry that the agent has failed to make; where the stakes of error are not high, withholding avoids no serious risk. Neither ground is present. The commitment is a voluntary policy and, by Principle~\ref{p:practical}, is permitted if its expected value is at least that of the alternatives. And it does not violate Principle~\ref{p:fit}, since a commitment is not a credence and does not require belief. So the agent who commits does nothing wrong.
\end{proof}

Theorem~\ref{t:james} covers the cases James cared about. In a genuine option that cannot be decided on intellectual grounds, evidence is not available by hypothesis, and the stakes of error are not high, since the option is forced. It also covers more modest cases: a scientist who adopts a working hypothesis that no available experiment can yet test, or a person who extends trust to a new colleague on whom there is no evidence either way.

The two theorems leave mixed cases to neither principle in its pure form. Where the stakes of error are high but evidence cannot be had, the agent should withhold judgment and act on her rational credence, which may well counsel not acting on $p$. Where evidence is available but the stakes are low, a belief on insufficient evidence violates Principle~\ref{p:fit} and is an epistemic defect, but it is not the wrong Clifford described: no one is exposed to a serious risk by it. This is why Clifford's principle, stated as holding \enquote{always, everywhere, and for anyone}, is false, and why \citet{vanInwagen1996} found it so hard to live by.

\section{The right to believe and uniqueness}\label{s:unique}

A permissivist reading of James is available: where the evidence cannot decide, several credences are rationally permitted, and the passions may choose among them. \citet{Kelly2014} develops a version of this thought from James's observation that the goals of believing truth and avoiding error can be weighted differently. EPI-17 argues that rationality is not permissive: on any body of evidence there is a unique rational credence in each proposition. We show that James's right is compatible with uniqueness.

\begin{proposition}[Reconciliation]\label{p:reconcile}
Where the evidence cannot decide whether $p$, the unique rational credence in $p$ lies strictly between the threshold for disbelief and the threshold for belief. James's right to believe, as stated in Definition~\ref{d:principles}, permits a commitment to $p$ while leaving that credence unchanged, and is consistent with Principle~\ref{p:fit}.
\end{proposition}

\begin{proof}
That the evidence cannot decide whether $p$ means that it is sufficient neither for $p$ nor for its negation; by Definition~\ref{d:sufficient}, the rational credence, which is unique by EPI-17, then lies strictly between the two thresholds. Principle~\ref{p:fit} requires that credence and forbids full belief. A commitment, by Definition~\ref{d:commit}, requires neither belief nor any particular credence, only the absence of disbelief, which the intermediate credence secures. So the agent may have the rational credence and the commitment together.
\end{proof}

Proposition~\ref{p:reconcile} is the position of REL-17-a, on which faith under underdetermination is a commitment that excludes disbelief without requiring belief, and on which evidentialism neither requires nor forbids it. It is also the form of the right that James's argument supports. His argument concerns options, which are choices of a course of life, and he himself characterized the maximum of liveness in a hypothesis as a willingness to act on it irrevocably. What the passional nature decides in a forced option is how to act, since some action is unavoidable. To read the right as a right to adopt a credence beyond the evidence would make it a right to do what, by EPI-11, cannot be done at will, and what, by EPI-17, is never rational.

James's further cases, in which faith in a fact helps to create it, fit the same account. If my commitment to a friendship makes it more likely that the friendship will flourish, then once I have committed, my evidence includes my commitment, and the rational credence in its success rises. The belief that may follow is then supported by evidence; it is not a credence chosen beyond the evidence. And since belief follows attention (EPI-11), a commitment that directs attention toward $p$ may in time bring evidence to light; any belief that results is assessed by that evidence, not by the commitment.

\section{Objections}

\begin{objection}[Momentous options are high-stakes]
James's options are momentous: much depends on them. To say that the stakes are not high in James's cases is to change the subject.
\end{objection}

\begin{reply}
Definition~\ref{d:stakes} concerns the stakes of error that withholding would avoid, and this is the notion Clifford's argument needs. The shipowner's error was wrong because he could have avoided the risk it imposed; a decision on which much depends but whose risks cannot be avoided by withholding is not one in which withholding is the responsible course. In a forced option, much is at stake, but none of it is protected by suspending judgment, and so the ground of Clifford's condemnation is absent.
\end{reply}

\begin{objection}[Clifford's argument from habit]
Clifford also argued that every belief on insufficient evidence, however trivial, weakens the habit of testing beliefs and contributes to a credulous society. This supports his principle everywhere, including where the stakes are low.
\end{objection}

\begin{reply}
The argument from habit supports the cultivation of a disposition to proportion belief to evidence, which Principle~\ref{p:fit} requires everywhere. It does not show that each trivial belief on insufficient evidence is a wrong for which the agent is answerable, any more than an argument for the habit of punctuality shows that every unpunctual act is a wrong. Clifford's principle, so understood, is an argument for an epistemic norm with universal scope and for a moral norm with the restricted scope of Theorem~\ref{t:clifford}.
\end{reply}

\begin{objection}[Wishful thinking]
The account licenses commitment to anything one wishes were true, provided the evidence is inconclusive.
\end{objection}

\begin{reply}
It licenses commitment only where evidence is not available, which excludes the many questions that could be settled by inquiry; only where withholding avoids no serious risk, which excludes most decisions with consequences for others; and only where the commitment is practically permitted, which requires that its expected value by the rational credence be at least that of the alternatives. Within those limits, it licenses no change in credence. Wishful thinking alters credence; the commitment of Theorem~\ref{t:james} does not.
\end{reply}

\begin{objection}[Thresholds may vary]
EPI-17 leaves open that the threshold for full belief varies with the agent's practical situation. If so, a Jamesian agent may lower the threshold and believe, not merely commit.
\end{objection}

\begin{reply}
Even if the threshold varied with the practical situation, it would not vary with the agent's choice, and by EPI-11 the agent could not adopt a belief at will in any case. The question whether practical considerations bear on the threshold for belief is the question of pragmatic encroachment, which our account does not need to settle. If they do, some of James's cases will be cases of permitted belief; the commitment that Theorem~\ref{t:james} permits is available either way.
\end{reply}

\section{Conclusion}

Clifford and James were each right about a domain. Where the stakes of error that withholding would avoid are high and further evidence can be had, believing on insufficient evidence results from a culpable failure of inquiry, and is wrong. Where the option is forced, so that withholding avoids nothing, and evidence cannot decide, Clifford's grounds lapse, and the agent may commit to a hypothesis on insufficient evidence without doing wrong. The right James defended is a right to commit, not a right to a credence the evidence does not support, and so it is consistent with the uniqueness of rational credence.

\begin{thebibliography}{99}\small
\bibitem[Clifford(1877)]{Clifford1877} Clifford, W.~K. (1877). The ethics of belief. \emph{Contemporary Review}, 29, 289--309.
\bibitem[Cohen(1992)]{Cohen1992} Cohen, L.~J. (1992). \emph{An Essay on Belief and Acceptance}. Clarendon Press.
\bibitem[Conee and Feldman(1985)]{ConeeFeldman1985} Conee, E. and Feldman, R. (1985). Evidentialism. \emph{Philosophical Studies}, 48(1), 15--34.
\bibitem[James(1896)]{James1896} James, W. (1896). The will to believe. \emph{The New World}, 5, 327--347.
\bibitem[Kelly(2014)]{Kelly2014} Kelly, T. (2014). Evidence can be permissive. In M.~Steup, J.~Turri and E.~Sosa (Eds.), \emph{Contemporary Debates in Epistemology} (2nd ed., pp.~298--312). Wiley-Blackwell.
\bibitem[Pace(2011)]{Pace2011} Pace, M. (2011). The epistemic value of moral considerations: Justification, moral encroachment, and James' \enquote{Will to Believe}. \emph{No\^us}, 45(2), 239--268.
\bibitem[van Inwagen(1996)]{vanInwagen1996} van Inwagen, P. (1996). It is wrong, everywhere, always, and for anyone, to believe anything upon insufficient evidence. In J. Jordan and D. Howard-Snyder (eds.), \emph{Faith, Freedom, and Rationality: Philosophy of Religion Today}. Rowman \& Littlefield.
\end{thebibliography}

\end{document}
